\documentclass[12pt]{article}

\usepackage[utf8]{inputenc}
\usepackage[T1]{fontenc}
\usepackage[polish]{babel}
\usepackage{listings}
\usepackage{xcolor}
\usepackage{geometry}
\usepackage{hyperref}
\usepackage{enumitem}

\geometry{a4paper, margin=2.5cm}

\lstset{
    language=Java,
    basicstyle=\ttfamily\small,
    keywordstyle=\color{blue},
    stringstyle=\color{red},
    commentstyle=\color{green!60!black},
    numbers=left,
    numberstyle=\tiny,
    frame=single,
    breaklines=true,
    showstringspaces=false
}

\title{Dokumentacja implementacyjna \\ \large Graficzny interfejs użytkownika -- wizualizacja grafów (Java Swing)}
\author{Mateusz Kamieniecki \\ Jan Rękas}
\date{\today}

\begin{document}

\maketitle
\tableofcontents
\newpage

% ─────────────────────────────────────────────────────────────────────────────
\section{Przegląd architektury}
% ─────────────────────────────────────────────────────────────────────────────

Projekt Java realizuje graficzny interfejs użytkownika do wizualizacji grafów,
współpracując z zewnętrznym silnikiem obliczeniowym napisanym w~języku C.
Całość opiera się na bibliotece \textbf{Java Swing} i~jest zorganizowana
w~pakiecie \texttt{graph}.

Architektura oprogramowania jest zbliżona do wzorca \textbf{MVC}
(Model--Widok--Kontroler):
\begin{itemize}
    \item \textbf{Model} -- klasy \texttt{Graph}, \texttt{Vertex}, \texttt{Edge}
          przechowują strukturę grafu bez żadnej logiki prezentacji.
    \item \textbf{Widok} -- \texttt{GraphPanel} (rysowanie grafu),
          \texttt{Footer} (pasek statusu), klasy panelów bocznych.
    \item \textbf{Kontroler} -- \texttt{GraphSettingsCard} i~\texttt{GraphInteractionPanel}
          reagują na akcje użytkownika, modyfikują model i~zlecają odświeżenie widoku.
\end{itemize}

Komunikacja z~silnikiem C odbywa się przez \texttt{ProcessBuilder} --
zewnętrzny plik \texttt{engine.exe} przyjmuje graf wejściowy i~zwraca
obliczone pozycje wierzchołków w~formacie tekstowym lub binarnym.

% ─────────────────────────────────────────────────────────────────────────────
\section{Klasy modelu danych}
% ─────────────────────────────────────────────────────────────────────────────

\subsection{graph.Vertex}

Klasa modelowa reprezentująca pojedynczy wierzchołek grafu.

\begin{lstlisting}
public class Vertex {
    public int    id;
    public double x, y;
}
\end{lstlisting}

\begin{description}[leftmargin=3cm, style=nextline]
    \item[\texttt{id}] Unikalny identyfikator całkowitoliczbowy wierzchołka.
    \item[\texttt{x}, \texttt{y}] Współrzędne ekranowe (w pikselach logicznych)
         wykorzystywane przez \texttt{GraphPanel} do rysowania i~przeciągania.
\end{description}

Współrzędne są modyfikowane bezpośrednio przez \texttt{GraphPanel} (podczas
przeciągania myszą) oraz przez \texttt{GraphSettingsCard} po wczytaniu pozycji
z~pliku. Klasa \emph{nie} przechowuje listy sąsiedztwa -- relacje między
wierzchołkami są modelowane wyłącznie przez krawędzie.

% ─────────────────────────────────────────────────────────────────────────────

\subsection{graph.Edge}

Klasa modelowa reprezentująca pojedynczą krawędź grafu.

\begin{lstlisting}
public class Edge {
    public String name;
    public Vertex from, to;
    public double weight;
}
\end{lstlisting}

\begin{description}[leftmargin=3cm, style=nextline]
    \item[\texttt{name}] Etykieta krawędzi wyświetlana na środku odcinka.
    \item[\texttt{from}, \texttt{to}] Referencje do wierzchołków końcowych (krawędź jest nieskierowana).
    \item[\texttt{weight}] Waga krawędzi wyświetlana obok nazwy.
\end{description}

Krawędź nie przechowuje własnych współrzędnych -- przy rysowaniu odczytuje je
bezpośrednio z~pól \texttt{x}, \texttt{y} obiektów \texttt{Vertex}.

% ─────────────────────────────────────────────────────────────────────────────

\subsection{graph.Graph}

Główna klasa modelu danych. Przechowuje listy wierzchołków i~krawędzi,
udostępnia metody do ich dodawania, usuwania i~wyszukiwania.

Kluczowe metody:
\begin{itemize}
    \item \texttt{addVertex(Vertex v)} -- dodaje wierzchołek do grafu.
    \item \texttt{addEdge(Edge e)} -- dodaje krawędź do grafu.
    \item \texttt{getVertexById(int id)} -- zwraca wierzchołek o~podanym
          identyfikatorze lub \texttt{null}.
    \item \texttt{getVertices()} / \texttt{getEdges()} -- zwracają listy
          (\texttt{ArrayList}) wierzchołków i~krawędzi.
    \item \texttt{clear()} -- usuwa wszystkie wierzchołki i~krawędzie.
\end{itemize}

Klasa jest czystym modelem danych -- nie zawiera logiki rysowania ani
algorytmów. Listy \texttt{ArrayList} są współdzielone referencją z
\texttt{GraphPanel}, więc każda modyfikacja modelu jest natychmiast widoczna
po wywołaniu \texttt{repaint()}.

% ─────────────────────────────────────────────────────────────────────────────
\section{Klasy widoku}
% ─────────────────────────────────────────────────────────────────────────────

\subsection{graph.Bar (klasa abstrakcyjna)}

Klasa bazowa dla wszystkich komponentów paskowych interfejsu (dolny pasek
statusu, potencjalny pasek narzędziowy). Dostarcza wspólnych stałych
stylistycznych:

\begin{lstlisting}
protected static final Color BG         = new Color(18, 18, 28);
protected static final Color BORDER_COL = new Color(40, 40, 60);
protected static final Color TEXT       = new Color(200, 200, 215);
protected static final Font  FONT       = new Font("Segoe UI", Font.PLAIN, 12);
\end{lstlisting}

Konstruktor ustawia tło i~czcionkę. Klasa nie zawiera logiki biznesowej --
służy wyłącznie ujednoliceniu wyglądu pasków.

% ─────────────────────────────────────────────────────────────────────────────

\subsection{graph.Footer (rozszerza Bar)}

Dolny pasek statusu informujący użytkownika o~bieżących akcjach
(ładowanie grafu, postęp algorytmu, błędy wejścia/wyjścia).

Kluczowa metoda:
\begin{lstlisting}
public void setStatus(String message) {
    SwingUtilities.invokeLater(() -> label.setText(message));
}
\end{lstlisting}

Użycie \texttt{SwingUtilities.invokeLater()} gwarantuje bezpieczeństwo
wątkowe -- metoda może być wywoływana z~dowolnego wątku, w~tym z~wątku
\texttt{SwingWorker} obsługującego długotrwałe obliczenia silnika.

Komponent zawiera \texttt{JLabel} z~domyślnym tekstem \emph{,,Gotowy''}.
Górna krawędź paska jest oddzielona linią o~kolorze \texttt{BORDER\_COL}.

% ─────────────────────────────────────────────────────────────────────────────

\subsection{graph.GraphPanel (rozszerza JPanel)}

Centralny komponent graficzny odpowiedzialny za wizualizację grafu
i~interakcję myszą. Rysowanie jest zaimplementowane samodzielnie przez
nadpisanie metody \texttt{paintComponent()}.

\subsubsection{Rysowanie}

\begin{lstlisting}
@Override
protected void paintComponent(Graphics g) {
    super.paintComponent(g);
    Graphics2D g2 = (Graphics2D) g;
    g2.setRenderingHint(RenderingHints.KEY_ANTIALIASING,
                        RenderingHints.VALUE_ANTIALIAS_ON);
    g2.translate(panX, panY);
    g2.scale(scale, scale);
    // rysowanie krawędzi, następnie wierzchołków
}
\end{lstlisting}

Transformacja \texttt{translate + scale} pozwala na płynne przesuwanie
i~skalowanie widoku bez modyfikacji współrzędnych modelu. Wierzchołki
rysowane są jako wypełnione okręgi (promień \texttt{RADIUS = 12} px),
krawędzie jako odcinki z~opcjonalnym tekstem pośrodku.

\subsubsection{Interakcja myszą}

Panel obsługuje trzy tryby interakcji:
\begin{itemize}
    \item \textbf{Przeciąganie wierzchołków} (lewy przycisk myszy) --
          po kliknięciu w~obszarze wierzchołka (\texttt{distance <= RADIUS})
          aktualizuje \texttt{v.x}, \texttt{v.y} w~modelu i~wywołuje
          \texttt{repaint()}.
    \item \textbf{Przesuwanie widoku (pan)} (środkowy przycisk myszy) --
          aktualizuje \texttt{panX}, \texttt{panY}.
    \item \textbf{Zoom} (kółko myszy) -- mnoży lub dzieli \texttt{scale}
          przez \texttt{1.1}.
\end{itemize}

\subsubsection{Metody publiczne}

\begin{description}[leftmargin=5cm, style=nextline]
    \item[\texttt{setZoom(double)}] Ustawia poziom powiększenia i~wywołuje
         \texttt{repaint()}.
    \item[\texttt{resetView()}] Zeruje przesunięcie widoku (\texttt{panX = panY = 0}).
    \item[\texttt{fitToView()}] Automatycznie dobiera zoom i~przesunięcie tak,
         by cały graf zmieścił się w~oknie z~80 px marginesem. Oblicza
         \texttt{minX}, \texttt{maxX}, \texttt{minY}, \texttt{maxY} i~wyznacza
         \texttt{scale = min(scaleX, scaleY)}, zachowując proporcje.
    \item[\texttt{setShowVertexDecorators(boolean)}] Pokazuje lub ukrywa
         identyfikatory wierzchołków.
    \item[\texttt{setShowEdgeDecorators(boolean)}] Pokazuje lub ukrywa
         etykiety i~wagi krawędzi.
    \item[\texttt{moveVertex(int id, double dx, double dy)}] Przesuwa wierzchołek
         o~podanym ID o~wektor $(\Delta x, \Delta y)$.
\end{description}

% ─────────────────────────────────────────────────────────────────────────────

\subsection{graph.LeftPanel}

Panel boczny działający jako kontener z~układem kartowym
(\texttt{CardLayout}). Zarządza przełączaniem między dwoma kartami:
\begin{itemize}
    \item \textbf{,,settings''} -- karta \texttt{GraphSettingsCard} (wybór
          pliku, trybu, algorytmu i~uruchomienie obliczania).
    \item \textbf{,,interaction''} -- karta \texttt{GraphInteractionPanel}
          (opcje wyświetlania, sterowanie widokiem, zapis pozycji).
\end{itemize}

Metody \texttt{showCard(String)} i~\texttt{getInteractionPanel()} umożliwiają
zewnętrznym komponentom przełączanie widoku.

% ─────────────────────────────────────────────────────────────────────────────
\section{Klasy kontrolera}
% ─────────────────────────────────────────────────────────────────────────────

\subsection{graph.GraphSettingsCard (rozszerza JPanel)}

Karta ustawień -- pierwszy ekran widoczny po uruchomieniu. Umożliwia:
\begin{enumerate}
    \item Wskazanie pliku wejściowego z~listą krawędzi grafu.
    \item Wybór trybu działania (algorytm lub wczytanie gotowych pozycji).
    \item Uruchomienie silnika obliczeniowego lub odczyt pozycji z~pliku.
\end{enumerate}

\subsubsection{Tryby pracy}

\begin{description}[leftmargin=5cm, style=nextline]
    \item[\texttt{radioFrucht}] Algorytm Fruchtermana-Reingolda (domyślny).
    \item[\texttt{radioTriang}] Algorytm triangulacji.
    \item[\texttt{radioLoadTxt}] Wczytanie pozycji z~pliku \texttt{.txt}.
    \item[\texttt{radioLoadBin}] Wczytanie pozycji z~pliku binarnego.
\end{description}

Wybór trybu dynamicznie pokazuje lub ukrywa sekcję \texttt{loadSection}
(pole ścieżki do pliku pozycji) i~\texttt{algoSection} (wybór formatu zapisu)
za pomocą \texttt{Runnable sync}.

\subsubsection{Metoda onExecute()}

Główny punkt wejścia po kliknięciu przycisku \emph{,,Wykonaj''}:
\begin{enumerate}
    \item Waliduje ścieżkę wejściową (niepusta, plik istnieje).
    \item Wczytuje krawędzie przez \texttt{GraphFileReader.readTextEdges()}.
    \item W~zależności od trybu:
    \begin{itemize}
        \item tryby \texttt{LoadTxt} / \texttt{LoadBin} -- wywołuje
              \texttt{applyPositionsTxt()} lub \texttt{applyPositionsBin()},
              następnie przełącza kartę;
        \item tryby algorytmiczne -- wywołuje \texttt{runEngine()}.
    \end{itemize}
\end{enumerate}

\subsubsection{Metoda runEngine()}

Uruchamia silnik C w~tle przy użyciu \texttt{SwingWorker}:

\begin{lstlisting}
SwingWorker<Void, String> worker = new SwingWorker<>() {
    @Override
    protected Void doInBackground() throws Exception {
        Process proc = new ProcessBuilder(cmd)
                .redirectErrorStream(true).start();
        // odczyt stdout i aktualizacja footer przez publish()
        int code = proc.waitFor();
        if (code != 0) throw new IOException(...);
        return null;
    }
    @Override
    protected void done() {
        // wczytanie wyników, przełączenie karty
    }
};
worker.execute();
\end{lstlisting}

Zastosowanie \texttt{SwingWorker} zapewnia, że interfejs nie blokuje się
podczas obliczeń. Metoda \texttt{process()} publikuje linie stdout silnika
do \texttt{Footer}, a~\texttt{done()} wczytuje wyniki i~przełącza na kartę
interakcji.

\subsubsection{Wyszukiwanie silnika -- findEnginePath()}

Metoda statyczna przeszukuje kilka kandydujących ścieżek:
\begin{lstlisting}
private static String findEnginePath() {
    String[] candidates = {
        "bin/engine/engine.exe",
        "app/bin/engine/engine.exe",
        "../bin/engine/engine.exe",
    };
    for (String c : candidates)
        if (new File(c).exists()) return new File(c).getAbsolutePath();
    return new File("bin/engine/engine.exe").getAbsolutePath();
}
\end{lstlisting}

\subsubsection{Metody pomocnicze (statyczne, dostępne dla innych klas)}

\begin{description}[leftmargin=4cm, style=nextline]
    \item[\texttt{styledButton(String)}] Tworzy przycisk z~ciemnym motywem.
    \item[\texttt{styledField()}] Tworzy pole tekstowe zgodne z~motywem.
    \item[\texttt{styledRadio(String)}] Tworzy radio button z~motywem.
    \item[\texttt{styledCheck(String)}] Tworzy checkbox z~motywem.
    \item[\texttt{sectionLabel(String)}] Tworzy szary nagłówek sekcji.
    \item[\texttt{box()}] Tworzy \texttt{JPanel} z~\texttt{BoxLayout.Y\_AXIS}.
    \item[\texttt{vgap(int)}] Tworzy pionowy odstęp (\texttt{Box.createVerticalStrut}).
    \item[\texttt{separator()}] Tworzy wizualną linię separatora.
    \item[\texttt{describe(Throwable)}] Zwraca czytelny komunikat wyjątku.
\end{description}

Metody te są \texttt{static}, dzięki czemu \texttt{GraphInteractionPanel}
importuje je statycznie i~zachowuje spójny wygląd bez duplikacji kodu.

% ─────────────────────────────────────────────────────────────────────────────

\subsection{graph.GraphInteractionPanel (rozszerza JPanel)}

Karta interakcji -- aktywna po załadowaniu grafu. Zawiera trzy sekcje:

\begin{description}[leftmargin=3.5cm, style=nextline]
    \item[Wyświetlanie] Dwa checkboxy: identyfikatory wierzchołków i~wagi
         krawędzi; każdy wywołuje odpowiednią metodę \texttt{setShow*()}
         na \texttt{GraphPanel}.
    \item[Widok] Trzy przyciski: \emph{Resetuj zoom}
         (\texttt{graphPanel.setZoom(1.0)}), \emph{Wyśrodkuj widok}
         (\texttt{graphPanel.resetView()}) oraz \emph{Dopasuj do widoku}
         (\texttt{graphPanel.fitToView()}).
    \item[Eksport] Przycisk \emph{Zapisz zmiany pozycji}, wywołujący
         metodę \texttt{savePositions()}.
\end{description}

\subsubsection{Metoda activate(String inputFilePath)}

Wywoływana przez \texttt{GraphSettingsCard} tuż przed przełączeniem karty.
Zapamiętuje ścieżkę pliku wejściowego, potrzebną do generowania nazwy
pliku wyjściowego.

\subsubsection{Metoda savePositions()}

Generuje nazwę pliku wyjściowego (\texttt{<baseName>\_newPositions.txt})
w~tym samym katalogu co plik wejściowy i~zapisuje aktualne współrzędne
wierzchołków przez \texttt{GraphFileReader.writeTextOutput()}.

% ─────────────────────────────────────────────────────────────────────────────
\section{Klasa graph.GraphFileReader}
% ─────────────────────────────────────────────────────────────────────────────

Klasa narzędziowa (wszystkie metody \texttt{static}) obsługująca wejście
i~wyjście plików grafu. Oddziela logikę parsowania od kontrolera.

\subsection{readTextEdges(String path)}

\textbf{Format wejściowy:} każda linia zawiera
\texttt{<nazwa> <fromId> <toId> <waga>} oddzielone białymi znakami.

\begin{lstlisting}
public static List<Edge> readTextEdges(String path) throws IOException
\end{lstlisting}

\textbf{Działanie:}
\begin{enumerate}
    \item Wczytuje plik linia po linii przez \texttt{BufferedReader}.
    \item Pomija puste linie.
    \item Waliduje, że każda linia ma co najmniej 4 tokeny; przy błędzie
          rzuca \texttt{IOException} z~numerem linii.
    \item Deduplikuje wierzchołki przez \texttt{LinkedHashMap<Integer, Vertex>}
          -- ten sam wierzchołek nie jest tworzony wielokrotnie, a~referencja
          jest współdzielona między krawędziami.
    \item Zwraca niepustą listę krawędzi; jeśli jest pusta -- rzuca
          \texttt{IOException}.
\end{enumerate}

\subsection{readTextVertices(String path)}

\textbf{Format wejściowy:} każda linia zawiera \texttt{<id> <x> <y>}.

\begin{lstlisting}
public static List<Vertex> readTextVertices(String path) throws IOException
\end{lstlisting}

Używany do wczytywania pozycji wierzchołków z~pliku \texttt{.txt} zapisanego
przez silnik C. Akceptuje zarówno kropkę, jak i~przecinek jako separator
dziesiętny (\texttt{replace(',', '.')}).

\subsection{writeTextOutput(Collection<Vertex> vertices, String path)}

\begin{lstlisting}
public static void writeTextOutput(Collection<Vertex> vertices, String path)
        throws IOException
\end{lstlisting}

Zapisuje pozycje wierzchołków do pliku tekstowego w~formacie
\texttt{\%d \%.2f \%.2f} (id, x, y) przez \texttt{PrintWriter}.

\subsection{readBinaryInput(String path)}

\begin{lstlisting}
public static List<Vertex> readBinaryInput(String path) throws IOException
\end{lstlisting}

Wczytuje plik binarny wygenerowany przez silnik C. Format rekordu (20 bajtów,
little-endian):

\begin{center}
\begin{tabular}{|c|c|c|}
    \hline
    \textbf{Offset} & \textbf{Typ} & \textbf{Znaczenie} \\
    \hline
    0 & \texttt{int} (4 B) & identyfikator wierzchołka \\
    4 & \texttt{double} (8 B) & współrzędna $x$ \\
    12 & \texttt{double} (8 B) & współrzędna $y$ \\
    \hline
\end{tabular}
\end{center}

Odczyt realizowany przez \texttt{ByteBuffer.wrap(allBytes).order(ByteOrder.LITTLE\_ENDIAN)}.
Pętla czyta kolejne rekordy dopóki bufor ma co najmniej 20 bajtów.

% ─────────────────────────────────────────────────────────────────────────────
\section{Klasy pomocnicze i kontener główny}
% ─────────────────────────────────────────────────────────────────────────────

\subsection{graph.MainFrame (rozszerza JFrame)}

Kontener główny okna aplikacji.

\textbf{Odpowiedzialność:}
\begin{itemize}
    \item Tworzenie instancji modelu (\texttt{Graph}) i~przykładowych danych
          (wierzchołki i~krawędzie do testów).
    \item Inicjalizacja widoków: \texttt{GraphPanel}, \texttt{Footer},
          \texttt{LeftPanel}.
    \item Ustawienie \texttt{BorderLayout}: \texttt{Footer} na dole,
          \texttt{LeftPanel} po lewej, \texttt{GraphPanel} na środku.
    \item Konfiguracja okna: tytuł, rozmiar, \texttt{EXIT\_ON\_CLOSE},
          wyśrodkowanie na ekranie.
\end{itemize}

\textbf{Uwaga:} Zgodnie z~wymaganiami należy dodać pasek menu
(\texttt{JMenuBar}) z~opcjami: wczytanie grafu z~pliku tekstowego i~binarnego,
zapis wyników oraz pomoc.

\subsection{graph.Main}

Klasa startowa. Metoda \texttt{main} uruchamia GUI w~wątku EDT:

\begin{lstlisting}
public static void main(String[] args) {
    SwingUtilities.invokeLater(() -> new MainFrame().setVisible(true));
}
\end{lstlisting}

Nie zawiera logiki biznesowej.

% ─────────────────────────────────────────────────────────────────────────────
\section{Wątkowość}
% ─────────────────────────────────────────────────────────────────────────────

Projekt przestrzega zasad bezpieczeństwa wątkowego Swing:

\begin{itemize}
    \item Wszystkie operacje na komponentach GUI odbywają się w~wątku EDT.
    \item \texttt{Footer.setStatus()} używa \texttt{SwingUtilities.invokeLater()}
          i~może być bezpiecznie wywoływana z~dowolnego wątku.
    \item Długotrwałe operacje (uruchomienie silnika C, oczekiwanie na wyniki)
          wykonywane są w~\texttt{SwingWorker.doInBackground()}, natomiast
          aktualizacja interfejsu -- w~\texttt{process()} i~\texttt{done()},
          które działają na EDT.
    \item Bezpośrednie modyfikacje modelu (\texttt{Vertex.x}, \texttt{Vertex.y})
          podczas przeciągania myszą odbywają się w~EDT (obsługa zdarzeń myszy).
\end{itemize}

% ─────────────────────────────────────────────────────────────────────────────
\section{Spójność komponentów i obsługa błędów}
% ─────────────────────────────────────────────────────────────────────────────

\begin{itemize}
    \item \texttt{GraphSettingsCard.onExecute()} waliduje ścieżki przed
          każdą operacją i~wyświetla \texttt{JOptionPane.ERROR\_MESSAGE}
          z~opisem błędu.
    \item Metoda \texttt{describe(Throwable)} (statyczna w~\texttt{GraphSettingsCard})
          wyodrębnia czytelny komunikat z~wyjątku, używana w~całym pakiecie.
    \item \texttt{GraphFileReader} rzuca \texttt{IOException} z~numerem linii
          przy błędzie parsowania, co pozwala użytkownikowi szybko znaleźć
          problem w~pliku wejściowym.
    \item Po pomyślnym wczytaniu grafu panel automatycznie przełącza się na
          kartę interakcji (\texttt{leftPanel.showCard("interaction")}),
          zapobiegając operowaniu na pustym widoku.
    \item Przycisk \emph{Wróć do ustawień} w~\texttt{GraphInteractionPanel}
          zawsze umożliwia powrót bez utraty wczytanego grafu.
\end{itemize}

% ─────────────────────────────────────────────────────────────────────────────
\section{Integracja z silnikiem C}
% ─────────────────────────────────────────────────────────────────────────────

Silnik obliczeniowy (\texttt{engine.exe}) jest osobnym procesem uruchamianym
przez \texttt{ProcessBuilder}. Interfejs komunikacji:

\begin{lstlisting}
engine.exe <sciezka_grafu> -a <algorytm> -o <sciezka_wyjsciowa> [-t|-b]
\end{lstlisting}

\begin{description}[leftmargin=3cm, style=nextline]
    \item[\texttt{-a}] Algorytm: \texttt{fruchterman} lub \texttt{triangulation}.
    \item[\texttt{-o}] Ścieżka bazowa pliku wyjściowego (bez rozszerzenia
         przy formacie binarnym).
    \item[\texttt{-t}] Zapis w~formacie tekstowym (plik \texttt{.txt}).
    \item[\texttt{-b}] Zapis w~formacie binarnym.
\end{description}

Silnik zapisuje wyniki w~podkatalogu \texttt{positions/} obok pliku grafu
wejściowego. Po zakończeniu procesu z~kodem 0 \texttt{GraphSettingsCard}
wczytuje wyniki i~stosuje je do modelu przez \texttt{applyPositionsTxt()}
lub \texttt{applyPositionsBin()}.

% ───────────────────────────────────────────────────────────────────────────

\end{document}